\SetPractica{Conservación de la materia y la energía}{Ley de la conservación de la materia y la energía}

\section{Introducción teórica}

Las bases de la química son la materia y la energía, y cualquier estudio de sus interacciones debe tener siempre en cuenta una propiedad fundamental de estas: la cantidad del conjunto de materia y energía en el universo se mantiene siempre constante. Esta propiedad realmente se divide en dos leyes científicas: la ley de la conservación de la materia y la ley de la conservación de la energía.

La primera suele expresarse en el siguiente postulado: ``La masa no se crea ni se destruye'' (Sterner, n.d.). A pesar de que tanto filosofías orientales como el jainismo y pensadores presocráticos ya habían explorado este concepto (Muni, 1983, p. 57) (Kirk and Raven, 1937, p. 291), este postulado no fue formulado formalmente hasta 1748 por el científico Mijaíl Lomonósov, sin embargo, no fue establecido como ley hasta 1785 por el químico Antoine Lavoisier, por esto  también se le conoce a esta ley como la ley de Lomonósov-Lavoisier (Significados, n.d.).

La segunda ley es la ley de la conservación de la energía, que se suele expresar en un postulado similar: ``La energía no se crea ni tampoco se destruye, sino que se transforma''. Esta ley asentó sus bases gracias a los experimentos de Julius von Mayer en 1842 (Munilla, 2023).

Ambas leyes nos hablan no sólo del carácter permanente de las cantidades de materia y energía en el universo, sino también de la posibilidad de ambas para transformarse. Una forma en la que esto se expresa es en los estados de agregación de la materia.

Las sustancias pueden existir, generalmente, en 3 estados; sólido, líquido y gaseoso. La diferencias entre cada uno de estos estados radica en el estado de los átomos que la componen, yendo de un estado altamente unido entre sí en los sólidos hasta un estado de alta separación entre las partículas en los gases, pasando por un estado intermedio en los líquidos. Una misma sustancia puede pasar de un estado de agregación a otro sin que ello implique que cambie su composición. El cambio de estado suele estar relacionado a la temperatura a la que la sustancia se exponga (Chang, 2002, p. 10-11).

Estos cambios de agregación, aunque cumplen con las leyes de conservación de materia y energía y no necesariamente cambian la composición de la sustancia, sí que pueden modificar sus propiedades físicas y químicas.

Una propiedad que se verá en esta práctica será la densidad, la cual se define como la relación entre la masa de un objeto y su volumen, y nos expresa la cantidad de sustancia que existe en un volumen dado (Densidad, n.d.).

Otra propiedad que se explorará es el punto de ebullición, el cual es la temperatura que un compuesto debe alcanzar para pasar del estado líquido al gaseoso. En el caso del agua este punto se encuentra en los 100C, pues es aquí que sus moléculas tienen la capacidad de separarse lo suficiente hasta cambiar de estado de agregación (Punto de ebullición, n.d.).

Otro concepto que se manejará es el valor de pH de diversas soluciones. Este valor representa la concentración de iones hidrógeno en una solución, y aunque existen diversas formas de medirlo, en esta práctica se utilizarán tiras que cambian a un color específico dependiendo el pH de la solución en que se mida. Los valores se representan en una escala que va del 0 al 14 y se divide en dos, del 0 al 6 para sustancias ácidas y del 8 al 14 para bases. El 7 representa una sustancia neutra (Lab, 2023).

Por último, se analizará el proceso de oxidación en frutas, el cual se da a causa de la absorción de radiación por parte de unas sustancias llamadas quinonas, que causan la coloración marrón en frutas oxidadas. Este proceso se puede ver afectado por la cantidad de luz, el calor y la cantidad de oxígeno en el ambiente en que se encuentre la fruta (La Oxidación De La Fruta, 2014).

\section{Objetivos}

\begin{itemize}
    \item Comprobar la ley de la conservación de la materia y la energía
    \item Identificar los estados de agregación de la materia y los cambios de fase.
    \item Determinar y distinguir propiedades físicas y químicas en diversas sustancias
\end{itemize}

\section{Materiales, equipos y reactivos}

\begin{table}[H]
\centering{\small\setstretch{1.25}
\begin{tabular}{|m{0.3\linewidth}|m{0.3\linewidth}|m{0.3\linewidth}|} \hline
    
    \thead{Equipos} & \thead{Materiales} & \thead{Reactivos} \\ \hline
    
    {Piseta \par}
    {Balanza digital \par}
    {Mechero Bunsen \par}
    {Soporte para mechero \par}
    {Malla de asbesto \par}
    {Termómetro \par}
    {Tiras de pH o potenciómetro \par}
    {Mortero y pistilo}
    
    &
    
    {Probeta \par}
    {Vasos de precipitado \par}
    {Varilla de agitación \par}
    {Matraz Erlenmeyer \par}
    {Vidrio de reloj \par}
    {Cuchillo \par}
    {Globo \par}
    
    & 
    
    {Sal \par}
    {Manzana \par}
    {Plátano \par}
    {Alka Seltzer\textsuperscript{\textregistered} \par}
    
    \\ \hline
\end{tabular}}\end{table}

\section{Procedimientos}

\subsection{Densidad}
\begin{enumerate}   
    \item Registrar el peso de los materiales de medición (probeta y vaso de pp).
    \item Medir 50mL en la de agua de agua, pesar y calcular la densidad.
    \item Vaciar el agua en el vaso de pp y agregar 1gr de NaCl (sal).
    \item Pesar de nuevo el vaso de precipitado y calcular la densidad de la mezcla.
\end{enumerate}

\subsection{Ebullición}
\begin{enumerate}
    \item Colocar 50mL de agua en un vaso de precipitado y registrar su temperatura.
    \item Colocar el vaso de pp a calentar y comenzar a cronometrar.
    \item Cuando comience a hervir, registrar la temperatura del agua, detener el cronómetro y retirar el vaso de pp de la fuente de calor.
\end{enumerate}

\subsection{pH}
\begin{enumerate}
    \item Colocar de 10 a 15mL de cada reactivo (agua, Coca-Cola y alcohol desnaturalizado) en un vaso de precipitado diferente cada uno.
    \item Sumergir el medidor de pH (tira de pH o potenciómetro) y registrar el valor.  
\end{enumerate}

\subsection{Oxidación}
\begin{enumerate}
    \item Cortar una manzana y un plátano, empezar a tomar evidencia fotográfica cada 5 minutos y comenzar a cronometrar.
    \item Registrar el tiempo en que aparece la primera mancha.
\end{enumerate}

\subsection{}
\begin{enumerate}
    \item Pulverizar una Alka-Seltzer en el mortero.
    \item Vaciar la pastilla ya pulverizada en un globo.
    \item Agregar 50mL de agua en un matraz Erlenmeyer.
    \item Sin que el contenido del globo se vacíe aún en el matraz, colocar la boca del globo en la del matraz.
    \item Pesar este conjunto de matraz y globo y registrar el peso.
    \item Voltear el globo para que su contenido se vacíe en el matraz y dejar ocurrir la reacción.
    \item Volver a pesar el conjunto de materiales y comparar con el peso previo a la reacción.
\end{enumerate}